\subsection{A reproducible direction-only localization response}

The robust result is a repeated head-level localization response, not a
universal classifier correction. In each of the RetinaMNIST CE, RetinaMNIST
RPS, Solar CE, and Solar RPS settings, all four confirmatory backbone seeds
had negative $\Delta\mathrm{MAE}$ under the frozen A-to-C comparison
(Table~\ref{tab:confirmatory-h1}). The direction of the effect therefore
replicated across the four tested dataset-objective settings. The replication
unit is the independently trained backbone seed. It is not an OOF fold, an
individual sample, or a separate dataset.

The A/C comparison constrains the interpretation more tightly than a full-head
retraining comparison. Condition C holds the representation, class-specific
row norms, and biases fixed while adapting classifier direction. The result
therefore supports the conclusion that the direction component of the
classifier head is a reproducible actionable component of rare-end localization
in the studied settings. It does not show that direction is the only relevant
head component, that all rare-end errors have the same source, or that the
intervention is universally preferable for prediction. In particular, the
localization response does not imply a global accuracy, calibration, risk, or
uncertainty-quantification improvement.

The descriptive confidence intervals in Table~\ref{tab:confirmatory-h1} also
illustrate why the frozen sign-consistency gate should be reported directly.
The RetinaMNIST CE interval crosses zero, but all four predeclared backbone
seeds improved under the paired endpoint. The setting consequently satisfies
the H1 criterion without requiring an additional significance decision. This
separation between a replication rule and descriptive uncertainty avoids
recasting a small number of backbone realizations as a large collection of
independent experiments.

\subsection{Endpoint-specific geometry is a setting-dependent diagnostic}

The direction-only effect is more robust than the centroid-based explanation of
that effect. H2a tests whether an endpoint-versus-adjacent representation margin
adds information about final C inward shrinkage beyond original-head state and
generic centroid separation. The result is strong in both Solar settings,
partial in RetinaMNIST CE, and not supported in RetinaMNIST RPS
(Table~\ref{tab:confirmatory-h2a}). Thus, endpoint-specific representation
geometry can provide explanatory value beyond the prespecified controls, but
it does not do so uniformly across the studied settings.

The contrast between Solar and Retina is informative. In RetinaMNIST CE, the
coefficient has the expected sign in all four seeds, but incremental
leave-one-out value is positive in only two. In RetinaMNIST RPS, only two seeds
have the expected coefficient sign. These outcomes do not support weakening the
Retina RPS negative result or replacing H2a with a more favorable analysis.
H2b remains secondary and cannot alter the H2a verdict. The appropriate
interpretation is that centroid geometry is a setting-dependent diagnostic of
post-adaptation localization, not a universal mechanism of recovery.

This boundary also separates two questions that can otherwise be conflated.
H1 asks whether direction-only adaptation improves rare-end localization. H2a
asks whether a particular representation diagnostic helps account for final
localization after that adaptation. A stable H1 response can therefore coexist
with a mixed H2a result. The present evidence supports the first question more
strongly than the second.

\subsection{Relation to long-tail classifier and representation studies}

The contribution is the ordinal localization object and the controlled
diagnostic, rather than classifier retraining or classifier geometry itself.
Balanced classifier retraining, decoupled representation/classifier training,
classifier norm and direction effects, and feature/classifier misalignment are
established themes in long-tail recognition
\cite{kang2020decoupling,alshammari2022weight,c2am2022angular,wang2026alignment}. This work does not
claim novelty for those components. It instead studies an ordinal failure mode:
a rare upper endpoint is routed toward adjacent interior classes, with a
corresponding shift in predictive location and inward shrinkage.

That distinction matters because nominal tail accuracy alone does not identify
where an error lies in the ordinal label space. The same exact-class error can
be adjacent to the endpoint or can be displaced several classes inward. The
frozen direction-only probe then tests whether part of this localization error
is head-actionable while the representation, row norms, and biases are held
fixed. This design provides a bounded classifier-head analysis of an
ordinal-specific endpoint phenomenon. It should not be presented as a new
balanced-head algorithm, a general result that classifier direction matters, or
a replacement for established long-tail methods.

\subsection{What the intervention establishes and what it does not}

The intervention shows that changing head direction alone can alter rare-end
localization. Historical seed-0 decomposition analyses give context for this
result: some rare-end samples are representation-inward under the
training-centroid diagnostic, whereas others retain rare-end-like geometry but
are mapped inward by the original head. The latter pattern is consistent with a
head-actionable component. These decomposition results are descriptive and
mechanism-forming; they are not a multi-seed causal classification of every
rare-end sample.

Several stronger interpretations are not supported. The analyses do not show
that the representation is irrelevant, that endpoint centroid position
causally determines recoverability, or that all representation-inward samples
are unrecoverable from frozen features. They also do not identify class
imbalance as the sole cause of inward localization. The support-severity study
is especially important on this point: lower class-4 support weakened direct
class-4 evidence, but localization outcomes were not cleanly monotonic. The
direction-only intervention is therefore best understood as an explanatory
probe, not as a production classifier selected for broad deployment.

\subsection{Controlled boundaries and limitations}

The controlled RetinaMNIST studies further bound the mechanism interpretation.
Within the frozen-RPS factorial, balanced sampling was the dominant observed
adaptation factor, while the CE-versus-RPS head objective had a smaller effect
within a sampling regime. This is a dataset-specific result from one frozen
backbone representation. It provides evidence about the training signal
associated with the beneficial direction change, but does not establish that
balanced sampling universally drives that change.

The support-severity grid provides a complementary boundary. Reducing class-4
support consistently degraded direct class evidence such as mean $p_4$ and the
class-4-versus-class-3 margin. Class-4 MAE, inward shrinkage, severe error, and
centroid-routing outcomes remained seed-variable and non-monotonic. Rare
support alone is therefore insufficient as an isolated causal account of
localization severity.

The confirmatory evidence is deliberately narrow. RetinaMNIST and Solar are
the two multi-seed confirmatory datasets, and Solar uses a fixed archived
readout rather than a new external cohort for each backbone seed. UTKFace has
provenance-clean seed-0 artifacts but lacks the seed-1--4 checkpoint set and is
historical supporting evidence only. Centroid distances are also a coarse
representation diagnostic, particularly when endpoint support differs across
datasets. These limitations do not remove the repeated H1 response, but they
define its scope and motivate a conservative mechanism claim: rare-end inward
localization has a reproducible direction-responsive head component in the
studied settings, while its geometry-based explanation is setting-dependent.
